
Fine-grained execution feedback improves reinforcement learning fine-tuning of code language models by providing token-level credit assignment, but this benefit depends on accurate error localization. This work demonstrates that Python exception types partition into categories with different traceback reliability: U\_line errors (SyntaxError, NameError, TypeError, etc.) achieve 100\% localization accuracy in our experiments, while U\_ignore errors (RuntimeError, RecursionError, MemoryError, etc.) achieve only 20\% accuracy. Applying token-level penalties at unreliable locations concentrates gradients at incorrect tokens---we measure significantly lower gradient concentration at ground-truth bug locations for U\_ignore errors (mean 1.398) compared to U\_line errors (mean 1.594), with p < $10^{-13}$. We propose error-type gating: apply fine-grained feedback only for reliably-localized errors, falling back to coarse feedback otherwise. In proof-of-concept experiments on APPS with CodeT5-small (60M parameters), gating shows a 4.61\% signal-to-noise ratio improvement (p = 0.112, direction confirmed but not statistically significant at $\alpha$ = 0.05) and enables the model to reach 30\% pass@1 when the ungated baseline does not within the same training budget. This work frames feedback granularity as a credit assignment reliability decision.

